\chapter{Discussion} \label{sec:discussion}

The FW-H code has been validated for monopole sources. The sensitivity studies confirmed that spatial discretisation affects the geometric accuracy of the spherical control surface. Similarly, temporal discretization must be fine enough to prevent the clipping of harmonic peaks for frequencies of interest.

The co-translating monopole case presented an interesting point. Despite having a zero relative Mach number ($M_r = 0$), the motion of the source through the quiescent fluid causes a phase shift of $\sim\pi/2$ compared to the stationary case. This is due to the proportionality of the convective term to the source $Q(\tau)$ rather than its derivative $\partial Q/\partial\tau$. Furthermore, the emission distance $r$ affects the retarded time calculation ($\tau = t - r/c$) as the wave moves between the translating points.

However, the dipole validation highlighted the limitations of explicitly deriving velocity terms. The results provide good agreement only when their assumptions are met. Assuming an analytical solution by a zero relative Mach number removes the $(1-M_r)$ factor, causing the analytical solution incapable of predicting moving sources. Similarly, the plane wave assumption fails with near-field control surface. The dual-monopole method alleviates these limitations by superimposing two physical monopoles, which preserves the convective terms without requiring explicit integration.

The limitations of the analytical boundary conditions are also shown in the plane wave test case (Figure \ref{fig:dual_far_valid}). The noticeably lower pressure amplitude  during the source's approach is an error due to imposing a static $x$-direction contribution on the dipole. As the source translates far downstream, the radiation vector creates an asymptote with the $x$-axis, such that the numerical and analytical results eventually match.

As briefly mentioned for the dual monopole results in sub-subchapter \ref{subsubsec:dipole_dual_monopole_validation}, the dual monopole approach shows numerical artefacts as the source moves away from the observer (Figure \ref{fig:dual_far_invalid}). This is likely due to a floating-point truncation error in the numerical code, as it lacks the required prevision for such small perturbation values. As most aeroacoustics codes solve in double-precision, the code needs to be updated to resolve these small perturbations for multiple sources.

To improve the accuracy of the dipole validation, the static $x$-direction constraint could be resolved by dynamically vectoring the dipole orientation normal to the control surface relative to the observer. Furthermore, the dipole source $f_i$ definition could be improved by using a non-stationary monopole equation to preserve the Doppler factor. Given these identified limitations in the analytical baseline, further studies to validate the dipoles are required before using the code to process CFD data.